\documentclass[UTF8,zihao=-4,scheme=chinese]{ctexart}

% Page and typography
\usepackage[
  a4paper,
  left=1.55cm,
  right=1.55cm,
  top=1.65cm,
  bottom=1.8cm,
  headheight=14pt,
  headsep=0.45cm,
  footskip=0.85cm
]{geometry}
\usepackage{setspace}
\setstretch{1.12}
\setlength{\parindent}{2em}
\setlength{\parskip}{0.25em}

% Math and diagrams
\usepackage{amsmath,amssymb,mathtools,bm}
\usepackage{tikz}
\usepackage{pgfplots}
\pgfplotsset{compat=1.18}

% Lists, tables, links
\usepackage{booktabs,array,tabularx}
\usepackage{enumitem}
\setlist[itemize]{leftmargin=2em,itemsep=0.15em,topsep=0.25em}
\setlist[enumerate]{leftmargin=2.2em,itemsep=0.15em,topsep=0.25em}
\usepackage[hidelinks]{hyperref}
\hypersetup{
  pdftitle={理工科课程复习讲义},
  pdfauthor={Course Alchemist}
}

% Header and footer
\usepackage{fancyhdr}
\pagestyle{fancy}
\fancyhf{}
\fancyhead[L]{\small 理工科课程复习讲义}
\fancyhead[R]{\small \leftmark}
\fancyfoot[C]{\small \thepage}
\renewcommand{\headrulewidth}{0.4pt}
\renewcommand{\footrulewidth}{0pt}

% Colors and boxed environments
\usepackage[most]{tcolorbox}
\tcbuselibrary{breakable,skins,theorems}

\definecolor{CAInk}{HTML}{1F2937}
\definecolor{CASubtle}{HTML}{F8FAFC}
\definecolor{CADef}{HTML}{2563EB}
\definecolor{CAThm}{HTML}{7C3AED}
\definecolor{CALem}{HTML}{0891B2}
\definecolor{CAHint}{HTML}{059669}
\definecolor{CAExam}{HTML}{D97706}
\definecolor{CAPitfall}{HTML}{DC2626}
\definecolor{CANote}{HTML}{475569}
\definecolor{CAMethod}{HTML}{0F766E}
\definecolor{CAExample}{HTML}{4F46E5}

\tcbset{
  ca-base/.style={
    enhanced,
    breakable,
    sharp corners,
    boxrule=0.6pt,
    leftrule=2.4pt,
    arc=1.2mm,
    left=1.1mm,
    right=1.1mm,
    top=0.9mm,
    bottom=0.9mm,
    before skip=0.65em,
    after skip=0.65em,
    fonttitle=\bfseries,
    coltitle=CAInk
  }
}

\newtcbtheorem[number within=section]{definition}{定义}{
  ca-base,
  colback=CADef!5,
  colframe=CADef,
  colbacktitle=CADef!12
}{def}

\newtcbtheorem[number within=section]{theorem}{定理}{
  ca-base,
  colback=CAThm!5,
  colframe=CAThm,
  colbacktitle=CAThm!12
}{thm}

\newtcbtheorem[number within=section]{lemma}{引理}{
  ca-base,
  colback=CALem!5,
  colframe=CALem,
  colbacktitle=CALem!12
}{lem}

\newtcbtheorem[number within=section]{hint}{提示}{
  ca-base,
  colback=CAHint!5,
  colframe=CAHint,
  colbacktitle=CAHint!12
}{hint}

\newtcbtheorem[number within=section]{examspot}{考点}{
  ca-base,
  colback=CAExam!7,
  colframe=CAExam,
  colbacktitle=CAExam!14
}{exam}

\newtcbtheorem[number within=section]{pitfall}{易错点}{
  ca-base,
  colback=CAPitfall!5,
  colframe=CAPitfall,
  colbacktitle=CAPitfall!12
}{pit}

\newtcbtheorem[number within=section]{attention}{关注}{
  ca-base,
  colback=CANote!6,
  colframe=CANote,
  colbacktitle=CANote!14
}{att}

\newtcbtheorem[number within=section]{method}{方法}{
  ca-base,
  colback=CAMethod!5,
  colframe=CAMethod,
  colbacktitle=CAMethod!12
}{mtd}

\newtcbtheorem[number within=section]{example}{例题}{
  ca-base,
  colback=CAExample!5,
  colframe=CAExample,
  colbacktitle=CAExample!12
}{ex}

\newtcolorbox{solutionbox}[1][]{%
  ca-base,
  title={解答},
  colback=CASubtle,
  colframe=CAInk!55,
  colbacktitle=CAInk!8,
  #1
}

% Common commands
\newcommand{\R}{\mathbb{R}}
\newcommand{\N}{\mathbb{N}}
\newcommand{\Z}{\mathbb{Z}}
\newcommand{\dd}{\,\mathrm{d}}
\newcommand{\CourseName}{课程名称}
\newcommand{\LectureTitle}{复习讲义标题}

\title{\bfseries \LectureTitle}
\author{Course Alchemist}
\date{\today}

\begin{document}

\maketitle
\tableofcontents

\section{模板示例}

这一页用于展示版面、边距和颜色框。实际生成讲义时，可以删除本节示例内容，只保留导言区和环境定义。

\begin{definition}{连续性}{continuity}
设函数 $f$ 在点 $x_0$ 的某邻域内有定义。若
\[
  \lim_{x\to x_0} f(x)=f(x_0),
\]
则称 $f$ 在 $x_0$ 处连续。
\end{definition}

\begin{theorem}{介值定理}{ivt}
若 $f$ 在闭区间 $[a,b]$ 上连续，且 $f(a)\neq f(b)$，则对任意介于 $f(a)$ 与 $f(b)$ 之间的数 $y$，存在 $\xi\in(a,b)$，使得 $f(\xi)=y$。
\end{theorem}

\begin{lemma}{符号保持性}{sign-preserving}
若 $f$ 在 $x_0$ 处连续且 $f(x_0)>0$，则存在 $\delta>0$，使得当 $|x-x_0|<\delta$ 时，有 $f(x)>0$。
\end{lemma}

\begin{examspot}{常考方式}{continuity-exam}
连续性常与零点存在性、极限计算、分段函数参数确定结合考查。遇到分段函数时，优先检查分界点处的左右极限与函数值。
\end{examspot}

\begin{hint}{做题提示}{continuity-hint}
证明存在性结论时，先尝试构造辅助函数，再检查它是否满足闭区间连续性和端点异号条件。
\end{hint}

\begin{pitfall}{忽略定理条件}{ivt-pitfall}
介值定理要求函数在整个闭区间上连续。只知道端点函数值异号，不足以推出区间内一定存在零点。
\end{pitfall}

\begin{attention}{需要关注的点}{attention-example}
讲义正文应尽量把定义、定理、例题和方法之间的逻辑关系写清楚，而不是只堆公式。复习时应关注“条件为何必要”和“题目如何触发该工具”。
\end{attention}

\begin{method}{分段函数连续性检查}{piecewise-method}
检查分段函数在分界点 $x_0$ 处连续，一般按三步进行：
\begin{enumerate}
  \item 计算左极限 $\lim_{x\to x_0^-}f(x)$；
  \item 计算右极限 $\lim_{x\to x_0^+}f(x)$；
  \item 比较左右极限与函数值 $f(x_0)$。
\end{enumerate}
\end{method}

\begin{example}{确定参数}{piecewise-example}
设
\[
f(x)=
\begin{cases}
ax+1, & x<1,\\
x^2+b, & x\ge 1.
\end{cases}
\]
若 $f$ 在 $x=1$ 处连续，求 $a,b$ 满足的关系。
\end{example}

\begin{solutionbox}
由左极限和右侧函数值相等可得
\[
  \lim_{x\to 1^-}(ax+1)=1+b,
\]
即 $a+1=1+b$，所以 $a=b$。因此连续性的条件为 $a=b$。
\end{solutionbox}

\end{document}
